\documentclass[UTF8,fontset=fandol,12pt,a4paper]{ctexart}
% Individual identity supplied by the author on 2026-11-09.
\newcommand{\AuthorName}{蒋佳豪}
\newcommand{\StudentID}{22671064}
\newcommand{\CourseName}{工程技术前沿}
\newcommand{\ReviewTitle}{面向科研任务的多智能体系统综述}
\newcommand{\ReviewSubtitle}{知识、记忆与协作的机制及证据边界}
\newcommand{\ReportDate}{2026年11月9日}

\usepackage{review}
\begin{document}
\reviewcover{文献阅读与写作准备（学习辅助，非提交稿）}

\section{要求核对与模板选择}
\subsection{这次需要交什么}
原始材料要求提交与课程相关的文献综述，包含问题引入、方法总结和自己的理解，
不少于3000字，图文并茂，文件为PDF，按“姓名+学号.pdf”命名。
课件截图给出的提交截止时间是2026年11月20日23:59；文字说明给出的互评截止日期为12月5日。
互评对象为同学，因此不宜把篇幅花在未经解释的框架名和内部术语上。

截图另明确禁止直接拷贝网络资料或直接提交AI生成内容。本材料只提供研究入口、
方法说明、比较问题和排版工具，不是满足要求的个人综述。
课堂概括本身不是论文原文，其中未经核对的人名、数值、机构或技术归因不用于书面结论。

\subsection{MP2模板哪些能用，哪些不能直接用}
可复用：ZJUI官方矢量标识、中文字体、衬线英文、A4页边距及较清楚的封面字号层级。
不直接复用：ECE2013课程名、教学班、四人组名录、实验章节、对局图表和默认英文提交方式。
本次改为个人中文文献综述，封面只保留课程、题目、作者、学号和日期；
正文改为问题、方法、比较、个人讨论和参考文献，采用12pt正文与编号引文。
这是个人维护的排版模板，不是学校统一指定的官方作业格式。

\subsection{建议的选题边界}
建议问题是：科研辅助系统如何在长任务与多角色协作中保持证据可追溯？
范围限定为基于大语言模型的科研辅助流程，重点比较知识检索、任务记忆和协作核验。
不把多智能体强化学习、群体机器人和所有AI for Science应用一并展开。
数据稀缺与生成增强可作展望中的边界问题，不额外引入一条与主线无关的完整综述。

课内概括提到“结构化知识增强、长上下文、记忆、自动科研以及验证瓶颈”，
因此该选题与课堂联系直接；但概括并未给出课堂方法的完整论文身份，
不能把GraphRAG论文的方法自动归为讲者具体工作的实现。

\clearpage
\section{概念地图：先区分三个问题}
\begin{figure}[htbp]\centering
\begin{tikzpicture}[node distance=9mm and 10mm,
 box/.style={draw=ZJUIBlue,rounded corners=2pt,align=center,text width=39mm,
 minimum height=17mm,font=\small},arr/.style={-{Stealth[length=2mm]},thick}]
 \node[box] (e) {外部证据\\文献、数据、引用};
 \node[box,right=of e] (m) {任务状态\\经历、失败、待办};
 \node[box,right=of m] (c) {协作控制\\分工、工具、反馈};
 \node[box,below=15mm of m,text width=69mm] (v) {核验与人的判断\\来源支持、实验复现、适用边界};
 \draw[arr] (e)--(v);\draw[arr] (m)--(v);\draw[arr] (c)--(v);
\end{tikzpicture}
\caption{用于组织阅读的概念图。为本材料归纳，不表示已实现或验证的新系统。}
\end{figure}

\begin{description}
\item[外部证据] 回答“结论从哪里来”。检索到相似文本不等于该文本真正支持当前断言；
应保存出处、上下文、版本及证据与结论的对应关系。
\item[任务状态] 回答“之前做过什么、哪些尝试失败”。它不等于事实数据库：
代理记录的推测应与实际观察分开，错误总结也不应因为进入长期记忆就被升级为事实。
\item[协作控制] 回答“由谁做、如何停止、怎样纠错”。多角色可以分工，
但相同模型、相同检索结果和相似提示也可能让它们共享同一个偏差。
\end{description}

上面关于证据与错误传播的区分是供作者检验的分析框架，不能写成所列八篇论文共同证明的结论。
如果作者赞同，应补充具体例子或论文证据；如果不赞同，可提出更合适的分类。

\section{文献检索记录}
检索日期：2026年11月9日。关键词包括：multi-agent conversation、AI research assistants、
GraphRAG、agent memory、long context和retrieval-augmented generation。
选取八篇公开原始论文，覆盖基础知识检索、上下文利用、行动反馈、协作和科研流程。
阅读卡依据论文摘要及方法、评估和局限部分的针对性核查，不宣称已逐字精读全部附录。
这是代表性阅读清单，不是系统综述，也不声称覆盖截至检索日的所有最新方法。
论文版本固定在参考文献中，全文下载校验记录保存在项目来源清单中。

\clearpage
\section{八篇原始论文的阅读卡}
\subsection{RAG：外部知识如何进入生成过程}
\noindent\textbf{来源与定位：}Lewis等，2020，所用版本v4；重点读第2节Models\cite{rag}。
\begin{itemize}
\item \textbf{文献事实：}模型结合预训练序列生成器与非参数检索记忆，
比较RAG-Sequence和RAG-Token两种对检索文档条件化的方式；原文包括训练与推断设计。
不能将其直接简化成今天任意“搜索网页后拼接提示”的全部实现细节。
\item \textbf{放在综述哪里：}作为知识增强的起点，用来解释模型参数知识与可检索文档的区别，
以及为什么文献更新不一定总要重新训练生成器。
\item \textbf{证据边界：}论文评估知识密集型NLP任务，不是多智能体科研系统的端到端评估；
不能用其问答结果直接推出科研结论可靠。
\item \textbf{作者需要回答：}一个引用存在，但引用内容不支持断言时，你认为错误发生在检索、
证据选择还是生成环节？请用自己读过的一条具体论断说明。
\end{itemize}

\subsection{GraphRAG：跨文档的全局问题}
\noindent\textbf{来源与定位：}Edge等，2024，所用版本v2；重点读索引与查询流程及第6节讨论\cite{graphrag}。
\begin{itemize}
\item \textbf{文献事实：}从语料抽取实体关系，形成分层社区摘要，再汇总面向查询的部分答案。
研究重点是面向语料整体的主题理解，而不只是检索一条局部事实。
\item \textbf{与课堂联系：}可对应“结构化知识”的线索，但原论文的社区总结流程，
不等于课内概括所述所有路径采样、结构约束或科学关系抽取细节。
\item \textbf{证据边界：}作者讨论了评估语料与问题类型的局限，并指出虚构率比较仍值得加强。
因此“回答更全面”不能自动改写为“所有事实更准确”。
\item \textbf{作者需要回答：}对“这篇文章用了什么数据集”和“这个领域有哪些研究路线”两类问题，
你会不会采用相同的检索结构？比较理由应落在任务差异上，而不是框架名上。
\end{itemize}

\clearpage
\subsection{Lost in the Middle：能放入不等于会利用}
\noindent\textbf{来源与定位：}Liu等，2023，所用版本v3；重点读多文档问答与键值检索实验\cite{middle}。
\begin{itemize}
\item \textbf{文献事实：}通过改变相关信息在输入中的位置及上下文长度，研究模型利用信息的表现。
所测条件下，中部相关信息往往比开头或结尾更难有效利用。
\item \textbf{可支撑的论述：}输入窗口容量和有效证据利用不是同一指标，
因此综述可以讨论压缩、检索、重排的必要性，而不是只比较窗口长度。
\item \textbf{证据边界：}研究针对当时的模型与构造任务；不能把其结果当成所有2026年模型
都具有同等退化幅度的证明，也不宜补写原文没有支持的统一百分比。
\item \textbf{作者需要回答：}如果你拿到几十篇论文，哪些内容应留在当前工作上下文，
哪些内容只需保留索引？给出一种具体取舍及可能遗漏的证据。
\end{itemize}

\subsection{ReAct：行动反馈与纯文本推理}
\noindent\textbf{来源与定位：}Yao等，2022，所用版本v3；重点读方法和不同任务的评估\cite{react}。
\begin{itemize}
\item \textbf{文献事实：}把推理轨迹与行动交替组织，使代理可通过外部环境获取反馈并调整计划；
原文使用问答、事实核查及交互任务评估。
\item \textbf{放在综述哪里：}作为“生成一个计划”到“根据观察继续行动”的桥梁，
解释为什么工具返回和执行结果比纯粹角色对话多提供了一种证据渠道。
\item \textbf{证据边界：}ReAct本身不是多代理协作框架；可观察的轨迹也不保证每一步推理正确，
工具错误、观察不足或检索失败仍需要讨论。
\item \textbf{作者需要回答：}代码能运行是否就证明科学假设正确？
请区分运行成功、指标复现、基线公平与结论成立四个层次。
\end{itemize}

\clearpage
\subsection{AutoGen：把协作写成可控制的对话}
\noindent\textbf{来源与定位：}Wu等，2023，所用版本v2；重点读第2节代理与对话编程\cite{autogen}。
\begin{itemize}
\item \textbf{文献事实：}采用可配置、可对话的代理，组合语言模型、工具及人的输入，
通过对话编程组织交互，支持多种任务。
\item \textbf{放在综述哪里：}用来说明分工、消息流、工具执行和人类参与如何成为系统设计的一部分。
重点不在列举角色名称，而在解释角色之间传递什么状态，以及何时结束。
\item \textbf{证据边界：}框架支持某种交互不等于它在所有任务上优于单代理。
原论文中的具体模型设置和示例应用，不能当作当前软件版本的操作说明。
\item \textbf{作者需要回答：}如果研究者和审稿者代理调用同一模型、看到同一材料，
二者一致是否构成独立核验？什么外部证据能够改变这一判断？
\end{itemize}

\subsection{Generative Agents：经历、检索、反思和规划}
\noindent\textbf{来源与定位：}Park等，2023，所用版本v2；重点读记忆架构与消融评估\cite{generative}。
\begin{itemize}
\item \textbf{文献事实：}记录经历，并按相关性、时间与重要性等因素检索，
将记忆反思和规划结合，用于交互式社会行为模拟。
\item \textbf{可迁移的概念：}科研任务也需要保存经历，但这只是机制类比。
作者可以讨论“做过什么”“为什么失败”“下一步计划”分别怎样记录。
\item \textbf{证据边界：}社会行为的可信模拟与科研事实的真实性不同。
消融支持其架构在原任务中的作用，不直接证明其可提升科研产出的正确性。
\item \textbf{作者需要回答：}把失败实验压缩成一句总结，可能丢失哪些条件？
你会保留哪些原始记录，以便之后推翻这句总结？
\end{itemize}

\clearpage
\subsection{MemGPT：记忆容量与上下文调度}
\noindent\textbf{来源与定位：}Packer等，2023，所用版本v2；重点读第2节内外部上下文设计\cite{memgpt}。
\begin{itemize}
\item \textbf{文献事实：}借鉴分层存储，区分当前主上下文与外部上下文，
通过显式调用移动信息，以支持超出有限窗口的文档和多轮交互任务。
\item \textbf{与其他方法比较：}它更关注状态调度，GraphRAG更关注语料组织；
二者不是天然互斥的整套替代品，但具体如何组合仍需系统设计和实验。
\item \textbf{证据边界：}记忆管理增加可访问的信息，不保证写入信息本身真实，
也不自动解决过期事实、冲突证据或代理之间的访问权限。
\item \textbf{作者需要回答：}当新论文否定一条旧结论时，是覆盖旧记忆，
还是保留两个版本及证据？说明你选择的更新规则及其代价。
\end{itemize}

\subsection{Agent Laboratory：科研流程与人的反馈}
\noindent\textbf{来源与定位：}Schmidgall等，2025，所用版本v2；重点读流程、人工参与评估及局限\cite{lab}。
\begin{itemize}
\item \textbf{文献事实：}从人的研究想法出发，组织文献阅读、实验和报告撰写流程，
允许人在各阶段提供反馈，评价包含研究产出及人的参与。
\item \textbf{放在综述哪里：}作为整合案例，把前面分散的知识、状态和协作问题放回科研链条，
观察错误在哪里被发现，或者在哪里可能被带入后续报告。
\item \textbf{证据边界：}论文报告的成本或性能收益都有特定比较条件。
阅读卡不搬运百分比作跨系统排名，也不将“能完成流程”视为普遍具有原创科学发现能力。
\item \textbf{作者需要回答：}在哪个阶段引入人的审查最值得？
请比较立题、实验设计、结果解释三个环节，而不只写“人机协同很重要”。
\end{itemize}

\clearpage
\section{横向比较：怎样从阅读变成综述}
\begin{table}[htbp]\centering\small
\caption{方法承担的不同功能。此表是阅读定位，不是性能排行榜。}
\begin{tabularx}{\linewidth}{@{}p{28mm}XX@{}}\toprule
类别与代表 & 主要问题 & 不能直接推出的结论\\\midrule
知识：RAG、GraphRAG & 怎样检索或组织外部材料 & 检索到就一定真实\\
状态：Generative Agents、MemGPT & 长任务中怎样保留并调用经历 & 记忆越多就越可靠\\
控制：ReAct、AutoGen & 怎样行动、分工及接收反馈 & 角色越多就越正确\\
流程：Agent Laboratory & 怎样贯通阅读、实验和写作 & 完成流程就证明创新\\\bottomrule
\end{tabularx}\end{table}

\subsection{比较时应固定什么}
若作者提出后续实验，应尽量固定输入材料、基础模型、检索范围、调用预算与人工反馈机会。
否则，“多代理比单代理效果好”可能同时包含更多调用、更强工具或更多人工帮助的影响。
若本作业没有实际运行系统，只能把这些作为评价建议，不写成实测结论。

\subsection{可以讨论哪些指标}
引用支持度关注断言与来源是否对应；复现完成度关注代码和数据能否重现结果；
任务成功率要先定义失败条件；成本可讨论调用与人的修订负担。
这些是建议的评价维度，不是八篇论文全部共同采用的指标。
不要在没有数据的情况下绘制准确率、成本改善或雷达评分图。

\subsection{适合互评读者的讲法}
先用一个熟悉的任务解释机制，例如“找论文—核对观点—写实验计划”，
再引入RAG、记忆和代理等术语。一个段落最好回答一个比较问题，
而不是连续介绍多个英文框架名。图表应帮助理解关系，不能替代正文论证或凑字数。

\clearpage
\section{正文规划与本人观点采集}
建议正文目标约3600字，分配为：问题550、方法1300、比较650、个人讨论850、结论250。
这一分配只是计划，最终必须核算实际完成的正文，提示文字和本阅读材料不计入作业字数。

请先写下三类原始笔记：一条你认可的结论及原文位置；一条你认为适用范围有限的结论；
一个你自己的课程或科研例子。每条尽量包含原因，而不仅是态度。
后续可以基于这些笔记整理段落、核对引文和改善表达，但不把本材料改名后直接提交。

\section{完成与提交检查}
\begin{enumerate}
\item 作者填写姓名与学号，并独立完成综述正文，阅读原文后补齐自己的判断。
\item 将排版样稿中的全部写作提示和待填写字段替换，摘要最后写。
\item 对每项文献事实核查引用，对个人推断明确给出理由，不编造实验或百分比。
\item 确认正文不少于3000字，且包括问题引入、方法总结与个人理解。
\item 检查至少一幅有解释作用的图和一张方法比较表，标明归纳或改绘来源。
\item 编译后逐页查看；正式PDF按“姓名+学号.pdf”命名，并在截止前提交。
\item 按课程安排完成互评。不要把截止日期等同于系统自动提交，本项目不自动登录教学平台。
\end{enumerate}

当前未完成项为：作者姓名、学号、本人原创正文和最终字数核验。
当前输出只在复核目录，不建立伪装成正式作业的提交文件。
参考文献均为真实公开论文；正文没有运行任何新的多智能体实验证明。

\clearpage
\begingroup\small\raggedright
\begin{thebibliography}{99}
\bibitem{rag} Lewis P, Perez E, Piktus A, et al. Retrieval-Augmented Generation for Knowledge-Intensive NLP Tasks[EB/OL]. 2020. arXiv:2005.11401v4. \url{https://arxiv.org/abs/2005.11401v4}.
\bibitem{graphrag} Edge D, Trinh H, Cheng N, et al. From Local to Global: A Graph RAG Approach to Query-Focused Summarization[EB/OL]. 2024. arXiv:2404.16130v2. \url{https://arxiv.org/abs/2404.16130v2}.
\bibitem{middle} Liu N F, Lin K, Hewitt J, et al. Lost in the Middle: How Language Models Use Long Contexts[EB/OL]. 2023. arXiv:2307.03172v3. \url{https://arxiv.org/abs/2307.03172v3}.
\bibitem{react} Yao S, Zhao J, Yu D, et al. ReAct: Synergizing Reasoning and Acting in Language Models[EB/OL]. 2022. arXiv:2210.03629v3. \url{https://arxiv.org/abs/2210.03629v3}.
\bibitem{autogen} Wu Q, Bansal G, Zhang J, et al. AutoGen: Enabling Next-Gen LLM Applications via Multi-Agent Conversation[EB/OL]. 2023. arXiv:2308.08155v2. \url{https://arxiv.org/abs/2308.08155v2}.
\bibitem{generative} Park J S, O'Brien J C, Cai C J, et al. Generative Agents: Interactive Simulacra of Human Behavior[EB/OL]. 2023. arXiv:2304.03442v2. \url{https://arxiv.org/abs/2304.03442v2}.
\bibitem{memgpt} Packer C, Wooders S, Lin K, et al. MemGPT: Towards LLMs as Operating Systems[EB/OL]. 2023. arXiv:2310.08560v2. \url{https://arxiv.org/abs/2310.08560v2}.
\bibitem{lab} Schmidgall S, Su Y, Wang Z, et al. Agent Laboratory: Using LLM Agents as Research Assistants[EB/OL]. 2025. arXiv:2501.04227v2. \url{https://arxiv.org/abs/2501.04227v2}.
\end{thebibliography}
\endgroup

\end{document}
